\chapter{The C toolbox}
\label{ch:c}

Operating-systems code is C (or, in SWEB, a careful subset of C++) with no
safety net: no garbage collector, no bounds checks, no exceptions. Most
bugs you will hunt in A1 are pointer and lifetime bugs that corrupt memory
silently and crash much later, somewhere else. This chapter is the minimum
toolkit: how to build with the right flags, how to let the compiler's
sanitizers find those bugs for you, and the three C ideas every later
chapter relies on --- pointers, lifetimes, and function pointers.

\section{Toolchain: from source to a running program}

\begin{lstlisting}[style=shell]
gcc -std=gnu17 -Wall -Wextra -g -O1 -fsanitize=address,undefined prog.c -o prog
\end{lstlisting}

\begin{description}
  \item[\code{-Wall -Wextra}] turn on the warnings that catch real bugs
    (unused results, sign mix-ups, missing returns). Treat a warning as a
    bug report you have not read yet. The course modules also turn a few
    warnings into errors, e.g.\ \code{-Werror=implicit-function-declaration}:
    calling a function without a prototype silently truncates pointers to
    \code{int} on 64-bit machines.
  \item[\code{-g}] debug information: line numbers in sanitizer reports
    and in \code{gdb}.
  \item[\code{-O0/-O1/-O2}] optimisation. Race bugs and missing
    \code{volatile}s often appear only at \code{-O2}, because the compiler
    then keeps values in registers and reorders memory accesses. The
    modules therefore test both a sanitizer build and an optimised build.
\end{description}

\code{make} rebuilds only what changed. Every module has a
\code{Makefile} with the targets \code{run} (examples) and \code{test}
(assignment). \code{make test IMPL=../../solutions/NN-...} builds the tests
against the reference solution instead of your own files.

\section{Sanitizers: let the compiler find your bugs}

A sanitizer is compiler instrumentation plus a run-time library. The
instrumented program checks every memory access (or every thread
interaction) and prints a report at the \emph{first} violation --- not at
the crash three functions later.

\begin{center}
\begin{tabular}{L{3.3cm}L{4.4cm}L{6.1cm}}
\toprule
\textbf{Sanitizer} & \textbf{Flag} & \textbf{Finds} \\
\midrule
AddressSanitizer & \code{-fsanitize=address} & out-of-bounds accesses (heap,
stack, globals), use after free, use after return, double free, leaks \\
Undefined\-Behavior\-Sanitizer & \code{-fsanitize=undefined} & signed overflow,
bad shifts, misaligned pointers, NULL dereference, out-of-range enum values \\
ThreadSanitizer & \code{-fsanitize=thread} & data races (chapter~3),
lock-order inversions (chapter~6) \\
\bottomrule
\end{tabular}
\end{center}

Reading a report: the first lines name the \emph{kind} of error
(\code{heap-use-after-free}, \code{stack-buffer-overflow}, \dots) and the
access; then come three stack traces --- where the bad access happened,
where the memory was allocated, and where it was freed. Start at the
first frame that is in your code.

\begin{warnbox}
ASan and TSan cannot be combined (they instrument memory differently),
and ThreadSanitizer fails on recent kernels with high address-space
randomisation (\emph{``unexpected memory mapping''}). The modules' test
runner starts TSan binaries under \code{setarch -R}, which disables
randomisation for that one process.
\end{warnbox}

\section{Pointers}

A pointer is an address plus a type. The type says how many bytes
\code{*p} reads and how far \code{p + 1} moves.

\begin{lstlisting}
int a[4] = {1, 2, 3, 4};
int *p = a;            /* an array decays to a pointer to its first element */
p + 2;                 /* address of a[2]: moves by 2 * sizeof(int) bytes    */
*(p + 2) == p[2];      /* identical: indexing is pointer arithmetic          */
sizeof a;              /* 16: the array; sizeof p would be 8: the pointer    */
\end{lstlisting}

\paragraph{Out-parameters.} C returns one value. A function that must
produce more writes through pointers the caller passes in ---
\code{pthread\_create(\&tid, ...)} stores the new thread's id in the
caller's variable, and every system call that reads data does the same
with a buffer.

\paragraph{\code{void *}.} A pointer to ``something''. It converts to and
from any object pointer without a cast in C, but cannot be dereferenced
or used in arithmetic until it is converted back to a real type. Generic
code (\code{qsort}, \code{pthread\_create}, a thread library's argument)
passes \code{void *} and trusts the caller to know the real type.

\paragraph{Integers as pointers.} Kernels constantly convert between
pointers and integers: a system call receives user pointers as plain
numbers in registers (SWEB: \code{size\_t}, or its \code{pointer} type).
\code{uintptr\_t}/\code{size\_t} are wide enough; \code{int} is not. This
is also why a \emph{number} chosen by a user must be validated before the
kernel treats it as an address (chapter~7).

\section{Lifetimes: stack, heap, static}

\begin{center}
\begin{tabular}{L{2.2cm}L{4.4cm}L{7.2cm}}
\toprule
\textbf{Storage} & \textbf{Lives} & \textbf{Typical bug} \\
\midrule
stack (locals) & until the function returns & returning or storing the
address of a local; a thread using a pointer to its creator's local after
the creator returned \\
heap (\code{malloc}) & until \code{free} & use after free, double free,
leak (never freed) \\
static (globals, \code{static}) & the whole program & shared between all
threads: needs synchronisation (chapter~3) \\
\bottomrule
\end{tabular}
\end{center}

The stack is the dangerous one for thread code: a loop that passes
\code{\&i} to every new thread passes the \emph{same} address, which keeps
changing --- and may no longer exist when the thread reads it
(\module{02}, \file{bug1\_loop\_index.c} and \file{bug2\_dead\_arguments.c}).

\begin{lstlisting}
int *broken(void) {
  int x = 42;
  return &x;        /* x dies here: the caller gets a dangling pointer */
}
\end{lstlisting}

AddressSanitizer reports this as \code{stack-use-after-return} (with
\code{detect\_stack\_use\_after\_return=1}, which the modules set).

\section{Function pointers and callbacks}

A function's name, used without parentheses, is its address. A function
pointer stores such an address and calls through it:

\begin{lstlisting}
int twice(int x) { return 2 * x; }
int (*f)(int) = twice;           /* pointer to function taking int, returning int */
f(21);                           /* 42 */
\end{lstlisting}

\paragraph{The callback pattern: function + context.} A library that
calls back into your code cannot know what data your code needs, so it
takes a function pointer \emph{and} an opaque \code{void *} that it hands
back unchanged:

\begin{lstlisting}
int pthread_create(pthread_t *thread, const pthread_attr_t *attr,
                   void *(*start_routine)(void *), void *arg);
\end{lstlisting}

\code{start\_routine} is ``what to run'', \code{arg} is ``with what''. The
library never looks inside \code{arg} --- which is also why the SWEB
kernel must \emph{not} validate it as a pointer (chapter~7): it may just
as well be the number 7 cast to \code{void *}. Returning a value works the
same way backwards: \code{start\_routine} returns a \code{void *} that
\code{pthread\_join} hands to the joiner.

\begin{swebbox}
SWEB is written in C++, but kernel C++: no exceptions, no standard
library (it brings its own \code{ustl} containers), fixed-width types
(\code{uint32}, \code{size\_t}, \code{pointer} for addresses). User
programs in \file{userspace/tests/} are plain C against SWEB's own small
libc (\file{userspace/libc}); the \code{pthread\_*} functions there are
empty stubs you will fill in.
\end{swebbox}

\begin{keybox}
\begin{itemize}
  \item Build with warnings on, and test with sanitizers \emph{and} with
    optimisation.
  \item Read a sanitizer report from the top: kind of error, bad access,
    allocation site, free site.
  \item Know where every object lives and how long; never hand a pointer
    to a local to code that outlives the function.
  \item Callbacks come as \emph{function pointer + opaque context}; the
    library never interprets the context.
\end{itemize}
\end{keybox}
